\section{M2 --- Pré-processamento}
\label{sec:m2}

Esta etapa transforma as séries de preço na representação consumida pelo modelo,
respeitando a ordem metodológica que evita vazamento temporal (Seção~\ref{sec:metodologia}).

\subsection{Log-retornos}
A feature principal é o retorno logarítmico diário $r_t = \ln(P_t/P_{t-1})$, calculado
sobre o fechamento ajustado. Por ser invariante a transformações multiplicativas de
escala, é imune a fatores de \emph{split}/grupamento já aplicados ao preço --- propriedade
relevante para o caso AMER3 (Seção~\ref{sec:m2-amer3}).

\subsection{Split temporal e normalização}
Aplica-se primeiro o \emph{split} temporal (treino até 2019; teste a partir de 2020) e só
então a normalização: o \texttt{MinMaxScaler} é ajustado \emph{exclusivamente} no treino e
aplicado a ambas as partições. Como consequência esperada, uma fração dos pontos de teste
extrapola o intervalo $[0,1]$ --- precisamente os extremos pós-2020 fora da faixa observada
no treino. Esse comportamento é desejável: é onde se concentram as anomalias, e o erro de
reconstrução tende a crescer ali. As janelas deslizantes ($\texttt{window\_size}=30$) são
geradas \emph{dentro} de cada partição, nunca cruzando a fronteira treino/teste. O
resultado são $\approx 2450$ janelas de treino e $\approx 1215$ de teste por ativo, no
formato $(n, 30, 1)$.

\subsection{Revisão da hipótese sobre o AMER3}
\label{sec:m2-amer3}
Em M1, levantou-se a hipótese de que o salto de log-retorno de $\approx +1{,}03$ em
14/11/2024 seria um \emph{artefato} de grupamento, e decidiu-se, provisoriamente,
``corrigir o fator''. A investigação conduzida em M2 \textbf{refutou essa hipótese} e
ilustra a importância de validar premissas com os dados:
\begin{itemize}[noitemsep]
  \item Os grupamentos reais do AMER3 ocorreram em \textbf{julho e agosto de 2024} (quatro
        eventos de razão 1:100 registrados pela fonte) e \textbf{já são corrigidos por}
        \texttt{auto\_adjust=True}: o log-retorno ajustado nessas datas é pequeno
        ($+0{,}09$; $+0{,}01$; $-0{,}18$; $+0{,}34$), com série contínua.
  \item Em 14/11/2024 o preço bruto saltou de R\$\,3{,}36 para R\$\,9{,}41 com volume
        cerca de três vezes a média e \emph{sem} \emph{split} registrado --- caracterizando
        um \textbf{movimento de mercado real} (especulação durante a recuperação judicial),
        assim como o salto de 15/08/2024.
  \item Como o log-retorno é invariante a escala, \emph{splits} já aplicados não produzem
        retorno espúrio.
\end{itemize}
\textbf{Decisão revista:} nenhuma correção é aplicada; o AMER3 segue o mesmo pipeline dos
demais ativos, e seus saltos de 2024 são tratados como anomalias reais (verdade narrativa
para M6). O registro completo está no ADR-0007.
